% Lecture 9, page 6 of 8 — English translation of the 1994 manuscript.
% Fragment intended for appending to courses/course-09.tex after page 5.
% Neither original Overleaf nor GitHub project has been modified.

\section*{Lecture 9 - Page 6}

\noindent\underline{Examples.}

\noindent 1) $G=S_n$, \qquad $A=\{1,2,\ldots,n\}$.
\[
S_n\times A\longrightarrow A
\]
\noindent $(\sigma,i)\longmapsto \sigma i=\sigma(i)\in A$, \qquad
$\sigma:\{1,2,\ldots,n\}\longrightarrow\{1,2,\ldots,n\}$.

\bigskip
\noindent We prove that this map is an action.

\bigskip
\noindent 2) $G$ an arbitrary group, $A=G$.\footnote{The two operation symbols are difficult to distinguish in the photograph. The notation $*$/$\cdot$ retains the manuscript's explicit distinction between external and internal operations.}
\[
\underset{\text{external op.}}{a*u}
\ \overset{\mathrm{def}}{=}\
\underset{\text{internal op.}}{a\cdot u}.
\]
% Editorial note: the manuscript's operation symbols are hard to
% distinguish; its external/internal operation distinction is retained.

\bigskip
\noindent 3) $G$ a group, $A=G$.
\[
a*u\ \overset{\mathrm{def}}{=}\ aua^{-1}
\qquad\text{(action by conjugation)}.
\]
\noindent We show that this is an action.

\bigskip
\noindent 4) $G=\left\{
\underbrace{\begin{pmatrix}
\cos\alpha & -\sin\alpha\\
\sin\alpha & \cos\alpha
\end{pmatrix}}_{R_\alpha}
\ \middle|\ \alpha\in\mathbb{R}\right\}$.

\noindent $R_\alpha R_\beta=R_{\alpha\beta}$.\footnote{The final subscript appears to read $\alpha\beta$; a ``$+$'' sign is not discernible in the photograph. The usual rotation identity is $R_\alpha R_\beta=R_{\alpha+\beta}$. The provisional manuscript reading is preserved here.}

\bigskip
\noindent If $\beta=-\alpha\ \Longrightarrow\
R_\alpha R_{-\alpha}=R_0=I_2=R_{-\alpha}R_\alpha$.

\bigskip
\noindent $G\leq GL_2(\mathbb{R})$.

\bigskip
\noindent Let $A=\mathbb{R}^2$ and write
$\begin{pmatrix}x\\y\end{pmatrix}\in\mathbb{R}^2$ (by notation).

\[
\begin{pmatrix}
\cos\alpha & -\sin\alpha\\
\sin\alpha & \cos\alpha
\end{pmatrix}
\begin{pmatrix}x\\y\end{pmatrix}
\ \overset{\mathrm{def}}{=}\
\begin{pmatrix}
x\cos\alpha-y\sin\alpha\\
x\sin\alpha+y\cos\alpha
\end{pmatrix}.
\]

\bigskip
\noindent\underline{Remark.} If $A$ is a $G$-set and $H\leq G$,
then $A$ is also an $H$-set:
\[
H\times A\longrightarrow A.
\]

\noindent If $\varphi:G\times A\longrightarrow A$, we define
$\varphi':H\times A\longrightarrow A$ by
\[
\varphi'(a,u)=\varphi(a,u).
\]
\noindent $\varphi'$ satisfies (1) and (2).

\bigskip
\noindent $\sigma\in S_n$ and $H=\langle\sigma\rangle\leq S_n$.
\[
H\times\{1,2,\ldots,n\}\longrightarrow\{1,2,\ldots,n\}
\]
\noindent $H=\{\sigma^h\mid h\in\mathbb{Z}\}$.
\[
\sigma^h\cdot i=\sigma^h(i)\in\{1,\ldots,n\}.
\]
